% 表示定理的几何：M=span{h_i}（蓝线）与正交补 M^perp（灰虚线）；读数与 f 相同的候选构成仿射集 P_M f + M^perp（绿线）；
% 其中离原点最近的是 P_M f，以原点为心、过 P_M f 的圆弧与该仿射集相切；f = P_M f + r，r 垂直于 M。
\begin{tikzpicture}[scale=1.15,note/.style={font=\FigureNoteFont\fontsize{7.5}{9}\selectfont,align=center}]
  \def\ang{20}
  \coordinate (O) at (0,0);
  \coordinate (u) at ({cos(\ang)},{sin(\ang)});
  \coordinate (v) at ({-sin(\ang)},{cos(\ang)});
  \coordinate (P) at ($2.2*(u)$);
  \coordinate (F) at ($(P)+1.25*(v)$);
  % arc of the circle through P centred at O
  \draw[FigureMuted,line width=.5pt,dash pattern=on 1.2pt off 1.2pt] ([shift={(O)}]-18:2.2) arc[start angle=-18,end angle=62,radius=2.2];
  % M and its orthogonal complement through O
  \draw[FigureBlue,line width=.9pt] ($-0.9*(u)$)--($3.4*(u)$) node[right,note,text=FigureBlue]{$M=\operatorname{span}\{h_1,\dots,h_n\}$};
  \draw[FigureMuted,line width=.6pt,dashed] ($-1.1*(v)$)--($1.9*(v)$) node[above,note,text=FigureMuted]{$M^{\perp}$};
  % affine set of equal readings
  \draw[FigureGreen,line width=.9pt] ($(P)-1.0*(v)$)--($(P)+2.0*(v)$) node[above,note,text=FigureGreen]{$\{g:Lg=Lf\}=P_Mf+M^{\perp}$};
  % right angle mark
  \draw[FigureStroke,line width=.45pt] ($(P)+0.16*(v)$)--($(P)+0.16*(v)-0.16*(u)$)--($(P)-0.16*(u)$);
  % vectors
  \draw[-{Stealth[length=4pt]},FigureStroke,line width=.7pt] (O)--(P);
  \draw[-{Stealth[length=4pt]},FigureRed,line width=.8pt] (P)--(F);
  \draw[-{Stealth[length=4pt]},FigureInk,line width=.6pt,dash pattern=on 2pt off 1.2pt] (O)--(F);
  \fill[FigureInk] (O) circle (1.1pt) node[below left=-1pt]{$0$};
  \fill[FigureBlue] (P) circle (1.4pt);
  \fill[FigureInk] (F) circle (1.4pt);
  \node[note,anchor=north west] at ($(P)+(0.04,-0.06)$) {$P_Mf$};
  \node[note,anchor=west] at ($(F)+(0.08,0)$) {$f$};
  \node[note,anchor=west,text=FigureRed] at ($(P)+0.62*(v)+(0.08,0)$) {$r$};
\end{tikzpicture}
